% archon:physics
% archon:covers ArchonPhysics/HarmonicModes.lean

\chapter{Finite harmonic normal modes}
\label{ch:ArchonPhysics_HarmonicModes}

For a finite positive-mass periodic lattice, this chapter records the
mass-weighted harmonic operator and its finite-dimensional normal-mode data.
All scalar occurrences below are named fixed-unit projections: a normal
coordinate $a$, its velocity projection $v$, and a frequency-square
$\omega^2$ are represented by \texttt{Real} only after choosing the project
units.  Thus they are not untyped stand-ins for a full physical state.  The
chapter proves finite spectral and oscillator facts only; it does not claim
Anderson localization, mixing, or thermalization.

\begin{definition}\leanok
[Periodic difference matrix]
\label{def:harmonic:difference-matrix}
\lean{ArchonPhysics.HarmonicModes.differenceMatrix}
\uses{def:lattice:site}
The matrix $D$ is the linear operator on site configurations satisfying
$(Dq)_i=q_{i+1}-q_i$.
\end{definition}
\begin{proof}

\leanok
The matrix is defined entrywise as the coefficient matrix of the periodic
forward-difference map.
\end{proof}

\begin{theorem}\leanok
[Difference matrix realizes the bond difference]
\label{thm:harmonic:difference-matrix-mulvec}
\lean{ArchonPhysics.HarmonicModes.differenceMatrix_mulVec}
\uses{def:harmonic:difference-matrix, def:lattice:forward-difference}
For every configuration $q$, matrix multiplication gives $Dq$ equal to the
locked periodic forward difference.
\end{theorem}
\begin{proof}

\leanok
At each site the row of $D$ has coefficient $-1$ at $i$ and $1$ at $i+1$.
The finite matrix-vector sum therefore reduces to $q_{i+1}-q_i$.
\end{proof}

\begin{definition}\leanok
[Mass-weighted difference]
\label{def:harmonic:mass-weighted-difference}
\lean{ArchonPhysics.HarmonicModes.massWeightedDifferenceMatrix}
\uses{def:harmonic:difference-matrix, def:lattice:mass-actions}
For a positive mass profile $m$, set $B=D M^{-1/2}$, where the diagonal
entry at site $i$ is $(\sqrt{m_i})^{-1}$.
\end{definition}
\begin{proof}

\leanok
Compose the periodic difference matrix with the diagonal inverse-square-root
mass action supplied by the lattice API.
\end{proof}

\begin{definition}\leanok
[Mass-weighted harmonic matrix]
\label{def:harmonic:mass-weighted-harmonic-matrix}
\lean{ArchonPhysics.HarmonicModes.massWeightedHarmonicMatrix}
\uses{def:harmonic:mass-weighted-difference}
Define the real symmetric dynamical matrix by $\Phi=B^{\mathsf T}B$.
\end{definition}
\begin{proof}

\leanok
This is the Gram matrix of the mass-weighted bond-difference operator.
\end{proof}

\begin{theorem}\leanok
[Gram factorization]
\label{thm:harmonic:gram-factorization}
\lean{ArchonPhysics.HarmonicModes.massWeightedHarmonicMatrix_eq_transpose_mul_self}
\uses{def:harmonic:mass-weighted-harmonic-matrix,
      def:harmonic:mass-weighted-difference}
The dynamical matrix is exactly $B^{\mathsf T}B$.
\end{theorem}
\begin{proof}

\leanok
Unfold the defining Gram-matrix expression.
\end{proof}

\begin{theorem}\leanok
[Nonnegative harmonic operator]
\label{thm:harmonic:matrix-pos-semidef}
\lean{ArchonPhysics.HarmonicModes.massWeightedHarmonicMatrix_posSemidef}
\uses{thm:harmonic:gram-factorization}
The matrix $\Phi$ is positive semidefinite.
\end{theorem}
\begin{proof}

\leanok
For every real vector $x$, $x^{\mathsf T}\Phi x=(Bx)^{\mathsf T}(Bx)$ is a
sum of squares and hence is nonnegative.
\end{proof}

\begin{definition}\leanok
[Fixed-unit modal energy]
\label{def:harmonic:modal-energy-api}
\lean{ArchonPhysics.HarmonicModes.modalEnergy}
For a fixed-unit frequency-square $\omega^2$, coordinate projection $a$, and
velocity projection $v$, define $E_{\omega^2}(a,v)=(v^2+\omega^2a^2)/2$.
\end{definition}
\begin{proof}

\leanok
This is the scalar energy of one projected harmonic degree of freedom in the
chosen mass and energy units.
\end{proof}

\begin{theorem}\leanok
[Conservation of fixed-unit modal energy]
\label{thm:harmonic:modal-energy-conserved-api}
\lean{ArchonPhysics.HarmonicModes.modalEnergy_conserved}
\uses{def:harmonic:modal-energy-api}
If $a'=v$ and $v'=-\omega^2a$, then $E_{\omega^2}(a(t),v(t))$ equals its value
at time zero for every $t$.
\end{theorem}
\begin{proof}

\leanok
Differentiate the numerator: $2vv'+2\omega^2aa'$
equals $2v(-\omega^2a)+2\omega^2av=0$.  A function with zero derivative on
$\mathbb R$ is constant, so evaluate the resulting identity at $0$.
\end{proof}

\begin{theorem}\leanok
[Nonnegative modal energy]
\label{thm:harmonic:modal-energy-nonneg-api}
\lean{ArchonPhysics.HarmonicModes.modalEnergy_nonneg}
\uses{def:harmonic:modal-energy-api}
If the frequency-square projection is nonnegative, then the modal energy is
nonnegative.
\end{theorem}
\begin{proof}

\leanok
Both $v^2$ and $\omega^2a^2$ are nonnegative; division by $2$ preserves the
order.
\end{proof}

\begin{definition}\leanok
[Mode frequency square]
\label{def:harmonic:mode-frequency-sq}
\lean{ArchonPhysics.HarmonicModes.modeFrequencySq}
\uses{def:harmonic:mass-weighted-harmonic-matrix}
For each normal-mode index $k$, $\omega_k^2$ is the eigenvalue of $\Phi$
associated with that mode.
\end{definition}
\begin{proof}

\leanok
Finite-dimensional real spectral data assign one eigenvalue to each vector of
an orthonormal eigenbasis.
\end{proof}

\begin{theorem}\leanok
[Nonnegative frequency squares]
\label{thm:harmonic:mode-frequency-sq-nonneg}
\lean{ArchonPhysics.HarmonicModes.modeFrequencySq_nonneg}
\uses{def:harmonic:mode-frequency-sq, thm:harmonic:matrix-pos-semidef}
Every mode frequency square is nonnegative.
\end{theorem}
\begin{proof}

\leanok
Apply positive semidefiniteness to a unit eigenvector.  Its quadratic form is
its eigenvalue, so that eigenvalue is nonnegative.
\end{proof}

\begin{definition}\leanok
[Normal-mode basis]
\label{def:harmonic:normal-mode-basis}
\lean{ArchonPhysics.HarmonicModes.normalModeBasis}
\uses{def:harmonic:mass-weighted-harmonic-matrix}
Choose an orthonormal real eigenbasis of the finite symmetric matrix $\Phi$.
\end{definition}
\begin{proof}

\leanok
The real finite-dimensional spectral theorem diagonalizes a symmetric matrix
by an orthonormal basis.
\end{proof}

\begin{theorem}\leanok
[Normal-mode eigen-equation]
\label{thm:harmonic:normal-mode-eigenvector}
\lean{ArchonPhysics.HarmonicModes.normalMode_eigenvector}
\uses{def:harmonic:normal-mode-basis, def:harmonic:mode-frequency-sq}
For every index $k$, $\Phi u^k=\omega_k^2u^k$.
\end{theorem}
\begin{proof}

\leanok
This is the eigenvector equation supplied by the selected spectral basis and
the corresponding eigenvalue definition.
\end{proof}

\begin{definition}\leanok
[Positive-frequency Physlib oscillator]
\label{def:harmonic:positive-mode-oscillator}
\lean{ArchonPhysics.HarmonicModes.positiveModeOscillator}
For a positive fixed-unit frequency $\omega$, construct the Physlib harmonic
oscillator with mass $1$ and spring constant $\omega^2$.
\end{definition}
\begin{proof}

\leanok
Instantiate the supplied oscillator structure with the stated positive
parameters.
\end{proof}

\begin{theorem}\leanok
[Positive-mode oscillator parameters]
\label{thm:harmonic:positive-mode-oscillator-spec}
\lean{ArchonPhysics.HarmonicModes.positiveModeOscillator_spec}
\uses{def:harmonic:positive-mode-oscillator}
The constructed oscillator has mass $1$ and spring constant $\omega^2$.
\end{theorem}
\begin{proof}

\leanok
Unfold the constructor and read off its two fields.
\end{proof}

\begin{theorem}\leanok
[Physlib positive-mode energy conservation]
\label{thm:harmonic:positive-mode-energy-conserved}
\lean{ArchonPhysics.HarmonicModes.positiveMode_energy_conserved}
\uses{def:harmonic:positive-mode-oscillator}
For a smooth trajectory satisfying that oscillator's equation of motion, its
Physlib energy at $t$ equals its energy at time zero.
\end{theorem}
\begin{proof}

\leanok
Apply the Physlib harmonic-oscillator energy-conservation theorem to the
constructed mass-one, positive-spring oscillator.
\end{proof}

\begin{definition}\leanok
[Translation mode]
\label{def:harmonic:translation-mode}
\lean{ArchonPhysics.HarmonicModes.translationMode}
\uses{def:lattice:mass}
Define the mass-weighted translation vector by $r_i=\sqrt{m_i}$.
\end{definition}
\begin{proof}

\leanok
This is the pointwise square-root mass profile, expressed in the fixed lattice
coordinate units.
\end{proof}

\begin{theorem}\leanok
[Translation zero mode]
\label{thm:harmonic:translation-mode-is-zero}
\lean{ArchonPhysics.HarmonicModes.translationMode_is_zero}
\uses{def:harmonic:translation-mode,
      def:harmonic:mass-weighted-harmonic-matrix}
The mass-weighted harmonic matrix sends the translation vector to zero.
\end{theorem}
\begin{proof}

\leanok
Multiplying $r$ by $M^{-1/2}$ gives the constant vector $1$.  Its periodic
forward difference is zero, hence $Br=0$ and $B^{\mathsf T}Br=0$.
\end{proof}

\begin{theorem}\leanok
[Translation mode is nonzero]
\label{thm:harmonic:translation-mode-ne-zero}
\lean{ArchonPhysics.HarmonicModes.translationMode_ne_zero}
\uses{def:harmonic:translation-mode, thm:lattice:mass-pos}
The translation vector is nonzero whenever the finite site type is nonempty.
\end{theorem}
\begin{proof}

\leanok
At any site, strict positivity of $m_i$ makes $\sqrt{m_i}$ strictly positive,
so the vector has a nonzero component.
\end{proof}

\begin{theorem}\leanok
[Harmonic data specification]
\label{thm:harmonic:data-spec}
\lean{ArchonPhysics.HarmonicModes.harmonicData_spec}
\uses{def:harmonic:difference-matrix,
      def:harmonic:mass-weighted-difference,
      def:harmonic:translation-mode,
      def:harmonic:modal-energy-api}
The mass-weighted difference has the stated diagonal factorization, the
translation vector has components $\sqrt{m_i}$, and the fixed-unit scalar
modal energy has its displayed formula.
\end{theorem}
\begin{proof}

\leanok
Unfold the three transparent definitions.  The matrix equality is the
definition of $DM^{-1/2}$, and the remaining two conjuncts are pointwise
definitional identities.
\end{proof}
